\section{Discussion and limitations}
\label{sec:discussion}
For a point query in an aligned scan pair, anatomy can order plausible translations cheaply, while a lesion scorer chooses among their positions.
Our equal-budget comparisons show that changing only \emph{where} the candidates lie, not how many are scored, improves matching in both cohorts.
Separating reach from selection explains the gain: diversification creates opportunities, and the scorer determines how many are realized.
The same accounting shows that a general anatomical embedding does not benefit from extra candidates that a lesion-specific scorer exploits.

\paragraph{Limitations.}
The task assumes an approximate current location and a known persistent lesion; clinical incident status and wrong-lesion associations require separate adjudication.
The query is simulated by jittering the annotated mark, and with the exact mark the screening advantage over the spline shrinks from 33 to 4 hits (Sec.~\ref{sec:query}); evaluating with queries from a real detector or real clicks is needed.
% TODO(group): agree on this wording.
A failed chest alignment can leave the lesion outside the local crop: in two staging pairs, 28 references lie more than 100\,mm outside the crop and every method errs by more than 200\,mm; these lesions remain in the denominators.
The screening cohort is retrospective, whereas the staging cohort is the independent test.
End-to-end time on native scans is dominated by the shared chest alignment.
Finally, we did not compare with learned registration models, which could provide a stronger continuous baseline.

\section{Conclusion}
\label{sec:conclusion}
We studied how to choose the candidate locations that a lesion scorer sees when matching a lung lesion across CT visits.
Ranking candidates by anatomical fit and keeping only those at least 2.1\,mm apart spreads a fixed budget over more distinct locations, and the complete list covers the full lattice to within that distance at about one-thirteenth of the cost.
On an independent staging cohort and a screening cohort, this simple rule outperforms a lesion-inclusive B-spline by about 10 points at 3\,mm while running about 19 times faster on the local step.
Our reach/selection decomposition shows that the gain comes from coverage, and that it pays off only with a scorer able to use the extra candidates.
Future work includes real query sources, scorers trained to exploit diverse candidate lists, and other anatomies where several nearby locations are plausible.
